\section{Does God Need Your Worship?}

Suppose the answer were yes. Praise would begin to resemble divine maintenance. Sacrifice or obedience would complete something missing in God, and the creature would sustain the Creator in the very act meant to acknowledge him. Worship would no longer tell the truth about either one.

Irenaeus rejects that picture directly. God did not form Adam because he needed a human being upon whom to depend. Before creation, the Word glorified the Father. Nor did God command service because he needed human obedience. Those who follow the light do not illuminate it; they are illuminated. Service profits God nothing, while the servant receives life, incorruption, and communion with God (\emph{Against Heresies} IV.14.1--2). The point governs Irenaeus's treatment of sacrifice as well. God does not hunger for oblations. He teaches a people to turn from idols, to practice justice, and to offer from a rightly ordered heart (IV.17.1; 18.1--6).

The result reverses the usual picture. We do not serve to supply a lack in God; we serve because we lack God. The act brings no gain to the One worshipped. It directs the worshipper toward communion with the good that no creature can produce for itself.

Aquinas asks what worship could possibly give to God. We offer not for his usefulness but for his glory and our good (ST II--II, q.~81, a.~6, ad 2). ``Glory'' here is not an increase in divine excellence. It is the creature's manifestation and acknowledgment of an excellence already infinite. Right worship creates no new truth about God; it brings the worshipper's act into conformity with truth.

\subsection{Can a free gift create a due?}

If God owes no creature the original gift of existence, how can a due arise at all? Aquinas distinguishes the orders carefully. God is not directed to creatures as to an end that perfects him. Within the order willed by divine wisdom, however, God gives each created nature what belongs to its condition; created dues are real but derivative (ST I, q.~21, a.~1, ad 3). The creature likewise possesses goods, powers, relations, and ends from which genuine obligations follow.

Creation begins with no antecedent claim in a nonexistent recipient. Once a rational creature exists, however, it is capable of acting fittingly or unfittingly toward the source from whom it exists. It can acknowledge truth or suppress it, honor the good or prefer a lesser good as ultimate, render justice or withhold it. The gratuity of the first gift does not leave moral life shapeless.

Two evasions now appear. One turns gift into contract: God gives in order to exact an equivalent. The other turns gift into moral indifference: because the gift was free, no response can be owed. A gift need not create commercial equivalence in order to establish a relation with moral form. The child's debt of honor to a parent, for example, is neither payment for conception nor a claim that parental goodness needs completion. The analogy is limited, but it shows why gratuity and due are not contraries.

The relation to God is more radical and less commensurable. Everything offered is already received: the body that kneels, the reason that judges, the will that consents, the bread and wine placed upon the altar, the hour reserved for worship. Nothing crosses from an independent estate into divine possession. The return cannot make the accounts equal because there are no such accounts.

Yet the creature's response is not morally empty. If the giver cannot be enriched and the return cannot be equal, what kind of justice remains?
